%TODO: 
% ARREGLAR FUNCION EXPONENCIAL CRECIENTE Y DECRECIENTE, DEPENDE DEL CONTEXTO
% MULETILLAS: ADEMAS, CABE DESTACAR
\documentclass{article}
\usepackage[utf8]{inputenc}
\usepackage[spanish]{babel}
\usepackage{graphicx, graphics, float, fancyhdr, titling}
\usepackage{listings}
\usepackage[a4paper, total={6in, 9.5in}]{geometry}
\usepackage{fancyhdr}
\usepackage{hyperref}   %para que funcione addcontentsline debe ser la ultima que se cargue

%\setcounter{secnumdepth}{-2}       %Poner solo esto si no se quieren numero delante de las secciones y niveles inferiores.

\renewcommand{\footrulewidth}{0.4pt}
\title{
\includegraphics[width=1.75in]{imagenes/UGR-Logo.png} \\
\vspace*{1in}
\textbf{Memoria de la práctica 2} \\
Animación por Ordenador \\
\vspace*{0.5in}}
\author{Andrés Merlo Trujillo \\
andresmerlo@correo.ugr.es \\
77147239H \\ 
\vspace*{0.5in} \\
E.T.S. de Ingenierías Informática y de Telecomunicación \\
\textbf{Universidad de Granada}} \date{\today}

\hypersetup{
    colorlinks=true,
    linkcolor=black,
    citecolor=black
}

\renewcommand\maketitlehooka{\null\mbox{}\vfill}
\renewcommand\maketitlehookd{\vfill\null}

\begin{document}
\begin{titlingpage}
\maketitle
\end{titlingpage}

\tableofcontents

\newpage

\pagestyle{fancy}   %a partir de comienza el header (se salta el indice y portada)
\fancyhead[L]{Andrés Merlo Trujillo}
\fancyhead[R]{Animación por Ordenador}
%\section{Ejercicio 1}
%\begin{figure}[H]
%    \centering
%    \includegraphics[width=\textwidth]{imagenes/passwdfile.png}
%    \vspace{10pt}
%    \footnotesize{Fuente: https://...}
%\end{figure}

\section{Ejercicio 1 - Pelota rodando}

En este ejercicio se pide implementar la animacion cuatro pelotas (o un haciendo un cuadrado) que parten desde un punto A y llegan al mismo tiempo a otro punto B haciendo, mediante curvas, animaciones distintas.

\bigskip

Cabe destacar que los \textit{keyframes} para todas las pelotas son exactamente los mismos, pero cambiando la forma de la curva. Entonces, los \textit{keyframes} son los siguientes:

\begin{enumerate}
    \item Instante 0: Posición inicial en el punto A.
    \item Instante 50: Posición final en el punto B.
\end{enumerate}

Y de manera gráfica, los \textit{keyframes} son:

%fotos de una bola nada mas con los keyframes

\bigskip

%foto con todas las pelotas xd

Estas animaciones las voy a dividir en subapartados para explicar mejor como son las curvas:

\subsection{Movimiento uniforme}

Para esta animacion es necesario usar un curva lineal, para que en cada instante de tiempo se avance siempre la misma distancia. La curva para esta pelota tiene la siguiente forma:

%foto de la curva

En la animación final se puede apreciar como mantiene la velocidad desde el primer momento hasta que acaba su recorrido.

\subsection{Aceleración}

Para realizar esta animación se debe usar una curva exponencial, cuya pendiente sea creciente. Esto va a hacer que progresivamente vaya recorriendo la pelota más distancia para la misma cantidad de tiempo. 

La curva tiene la siguiente forma:

%foto de la grafica

En la animación se puede apreciar como poco a poco la pelota va acelerando cada vez más, hasta para bruscamente en en el putno B.

Cabe decir que he modificado la forma de la curva para que se pueda apreciar mejor la aceleración en la animación resultante, ya que con la curva por defecto del programa no se puede apreciar del todo bien.

\subsection{Deceleración}

En este caso se debe usar una curva similar para la de aceleración, pero con una pendiente que vaya siendo cada vez menor. Esto resultará en que la pelota cada vez va a recorrer menos distancia para la misma cantidad de tiempo, haciendo que al final frene.

La curva tiene la siguiente forma:

%foto de la curva

En la aniamción resultante la pelota comienza moviéndose muy rápido y conforme se acerca a la posición final va frenando.

Al igual que con la aceleración, he modificado la curva para que el frenado sea más pronunciado.

\subsection{Aceleración y deceleración}

Esta animación es la combinación de la versión de aceleración y deceleración, haciendo que la pelota acelere cuando salga del punto A y vaya frenando cuando se acerque al punto B. 

La curva resultante es la unión de dichas versiones, haciendo que la pendiente sea progresivamente crenciente, hasta llegar a un punto en el que la pendiente decrece. 

Dicha curva es la siguiente:

%imagen de la curva

El resultado es que la pelota acelera al principio y cuando se acerca al punto final comienza a frenar. Ademas, como en los anteriroes subapartados, he modificado la curva para que el resultado sea mas facil de ver.


\section{Ejercicio 2 - Salto de coche}

En este ejercicio se pide animar un coche que arranca desde una posición, para luego subir una rampa y saltar por los aires, para finalmente cuando toca el suelo frenar.

% no se si poner como he hecho la rampa.

%reescribir esto
Cabe destacar que para animar la rotacion para subir a la rampa, he cambiado el pivote para poner justo en el eje de las ruedas traseras, ya que es como realmente se deberia mover al intentar subir dicha rampa. 

Además, la jerarquía utilizada ha sido la de dejar como padre el cubo más grande del vehículo y como hijos las ruedas y el cubo más pequeño que hace de pilares y cristales.


La animación consiste en 8 \textit{keyframes}, que son los siguientes:

\begin{enumerate}
    \item Instante 0: El vehículo se encuentra en la posición inicial.
    \item Instante 13: El vehículo ha realizado la aceleración.
    \item Instante 26: El vehículo se encuentra en el instante anterior de subir la rampa.
    \item Instante 28: El vehículo ha subido la rampa y se encuentra al principio de la misma.
    \item Instante 32: El vehículo se encuentra en el instante anterior de estar en el aire; es decir, se encuentra tocando el final de la rampa solo con la rueda trasera.
    \item Instante 44: El vehículo se encuentra en el punto más alto del lanzamiento y ya recto.
    \item Instante 56: El vehículo ha tocado el suelo después del salto.
    \item Instante 85: Finalmente, el vehículo frena hasta pararse del todo.
\end{enumerate}

A modo visual, los \textit{keyframes} son los siguientes:

%fotos de los keyframes

Y las curvas utilizadas para la animación:

%fotos de las curvas, separadas por partes (si no son muchas)

% reescribir lo de abajo con chatgpt
Como se puede ver en la curva roja (posición en el eje X), al principio la curva tiene forma exponencial creciente, para simula la aceleración. A continuación, todo el trayecto se mantiene lineal, hasta llegar a la parte de frenado, donde tiene forma exponencial para que frene y se pare.

En cuanto a la curva de color azul (Posición en el eje Z), solo es usada para animar la subida de la rampa y la trayectoria que sigue el vehículo cuando sale despedido de la rampa. Entonces, la parte en la que sube a la rampa es lineal, con la misma pendiente que la rampa. Mientras que la animación de la trayectoria se hace con dos funciones exponenciales, para animar la fuerza que realiza la gravedad en la parte de la subida y para animar la aceleración de la gravedad en la bajada. Cabe destacar que, en este caso, la forma que tenga la curva será reflejada en la animación directamente; es decir, si la trayectoria se hubiera hecho lineal, el lanzamiento hubiese tenido forma triangular.

En cuanto a la curva verde (Rotación en el eje Y), es usada para rotar el vehículo en el cambio de pendiente de la subida de la rampa y cuando se endereza durante el lanzamiento del vehículo.
% fin rescritura

\section{Ejercicio 3 - Pelota botando}

En este ejercicio se pide animar una pelota que rebota varias veces sobre una mesa y que finalmente se cae al suelo. 

La mesa la he realizado usando 5 cubos: uno grande para el tablero y otros 4 del mismo tamaño para hacer las patas. Además, la jerarquía utilizada ha sido de dejar como padre al tablero y como hijos las patas.

%rescribir
En cuanto a la animacion, he utilizado un factor de 2 para la altura del rebote y para la distancia recorrida entre rebotes; es decir, que en cada salto su altura y distancia de divide a la mitad. La única parte en la que no se sigue esto es en el primer lanzamiento horizontal, ya que probando a seguir este factor el resultado era menos realista, pero la altura sí lo respeta, como pedía el ejercicio. Esta aniamción la he realizado usando 10 \textit{keyframes}, que son:

\begin{enumerate}
    \item Instante 0: La pelota se encuentra en el aire, lista para ser lanzada a la mesa.
    \item Instante 15: La pelota ha tocado la mesa por primera vez.
    \item Instantes 25, 42, 54: La pelota ha rebotado y se encuentra en el aire.
    \item Instantes 35, 49, 59: La pelota ha tocado de nuevo el tablero de la mesa.
    \item Instante 71: La pelota se encuentra en el borde de la mesa, justo antes de que se caiga al suelo.
    \item Instante 84: La pelota ha caido y se encuentra en el suelo.
\end{enumerate}

% imagenes con los keyframes

Y para animar los rebotes y el movimiento de manera realista, se han usado las siguientes curvas en la pelota:

% imagenes separadas de las curvas

La curva de color rojo (Posición en el eje X) sigue una forma lineal para todos los \textit{keyframes}. También se puede observar como el conjunto de todos los \textit{keyframes}, al seguir el factor 2, tiene una forma exponenecial, haciendo que en la animación la pelota vaya recorriendo cada vez menos distancia. 
Esto se podría haber hecho directamente con dos \textit{keyframes} y dicha curva, pero se perdería demasiada precisión. % quizas quitar esto

Mientras que la curva de color azul (Posición en el eje Z) sigue las trayectorias que hace la bola. Para que sean realistas los botes, he utilizado una curva de desaceleración en la subida y otro de aceleración en la bajada, para simular el efecto de la gravedad en la pelota. Cuando cae de la mesa se sigue la misma curva de aceleración.

\section{Ejercicio 4 - Montacargas}

En este ejercicio se pide animar un montacargas que se mueve desde la izquierda hasta la derecha y para burscamente al final aplicando el principio de \textit{Overlap \& Follow Through}. Por tanto, la cadena del montacargas debe ser ``arrastrada'' por la plataforma y cuando esta pare, la cadena debe dar un latigazo en el sentido contrario, para luego detenerse en el sitio final.

%jerarquia de las piezas (reescribir)
La figura del montacargas la he realizado usando como plataforma un cubo achatado y para la cadena he usado como figuras básicas un cilindro y una esfera para la articulación. Usando estas figuras, he anidado las distintas secciones del brazo, haciendo que el hijo sea la pieza de abajo y el padre la de arriba. Así se consigue que el hijo herede las transofrmaciones del padre.

%foto de la jerarquia


En cuanto a la animacion, ha sido necesario utilizar 2 \textit{keyframes} para la plataforma y 6 \textit{keyframes} para cada una de las partes del brazo, pero que realizan exactamente lo mismo. 

\textit{Keframes} de la plataforma:


\begin{enumerate}
    \item Instante 0: Posición inicial, a la izquierda del todo.
    \item Instante 30: Posición final, a la derecha del todo.
\end{enumerate}

\textit{Keframes} generales del brazo:

\begin{enumerate}
    \item Instante 0: Se encuentra el segmento en reposo; es decir, en vertical.
    \item Instante 18: Se encuentra cada segmento del brazo rotado al máximo, siendo arrastrado por la plataforma en movimiento.
    \item Instante 30: Exactamente igual que antes, pero en este caso la plataforma ya ha parado.
    % arreglar parada en seco con otras palabras
    \item Instante 50: Cada segmento del brazo ahora se encuentra rotado al lado contrario, fruto de la parada en seco de la plataforma, haciendo que de un latigazo.
    \item Instante 64: Cada segmente vuelve a estar rotado hacia el otro lado, pero esta vez con una rotación algo menor que la primera vez, simulando las distintas fuerzas que actuan.
    \item Instante 76: Los segementos se encuentran de nuevo verticales, como en el instante 0.
\end{enumerate}
%keyframes

A modo visual, los \textit{keyframes} son los siguientes:

% explicar el significado del instante 18 (rescribir)
El instante 18 lo he puesto porque sin él la animación de rotar el brazo por la velocidad no se acababa hasta el punto en el que la plataforma parase, haciendo que la animación no quedase demasiado convincente. Haciendo esto, he conseguido que el brazo rote del todo antes de parar.

%curva
La curva para la aniamción de la plataforma es la siguiente:

%foto de la curva

Dicha curva lo que hace es, al principio acelera y cuando llega al punto final se para bruscamente (mediante la curva lineal). Esto lo he hecho asi para que el latigazo que realiza el brazo sea mas realista, ya que si se va frenando progresivamente no queda del todo bien.

Las curvas de los segmentos del brazo son muy similares entre si, por lo que solo voy a mostrar uno.

%foto de la curva

%rescribir
Esta curva lo que hace es seguir la aceleracion de la plataforma al principio, hasta llegar al instante 18 que he hablado anteriormente, asi se consigue que el movimiento sea mas natural. Despues se mantiene constante al no requerir rotacion ninguna. Cuando la plataforma para, se sigue una curva de desaceleracion, ya que la energia que tenia la pataforma se transfiere al brazo. Finalmente, se utilizan curvas de aceleración, desaceleracion y \textit{Slow in/Slow out} (union de ambas) para simular el baiben que hace hasta parar por las distintas fuerzas.


\end{document}
